\subsection{No compatible Picard variety over dual numbers}
\paragraph{D4a (Explicit counterexample to Artinian compatibility).}
Let $k=\overline{\mathbb F}_2$, $R=k[\epsilon]/(\epsilon^2)$,
and take two labelled copies $E_0,F_0$ of the ordinary elliptic
curve
\[
 y^2+xy=x^3+1,
 \qquad P=(0,1)\in F_0[2](k).
\]
Its discriminant is one and the coefficient of $xy$ is one, so it
is smooth ordinary; inversion sends $(x,y)$ to $(x,y+x)$,
showing that the displayed nonzero point has order two.
Set $A_0=E_0\times F_0$ and
$X_0=A_0/\langle(-1,t_P)\rangle$, as in D4o.
Choose the marked abelian deformation $A/R$ with Katz coordinates
\[
 (q_{ij})=\begin{pmatrix}1&1+\epsilon\\1&1\end{pmatrix}.
                                                        \tag{D4a.1}
\]
The involution and its translation lift: in D4o, equations
(I1)--(I3), choose $z_1=z_2=1$. Indeed
$(1+\epsilon)^2=1$ and the second row is $(1,1)$.
This gives an actual smooth projective quotient deformation
$X/R$ and a finite étale double cover $\pi:A\to X$.
Its torsion Picard scheme is flat, by D4o.

Nevertheless there is \emph{no} abelian $R$-scheme $B$ with a
homomorphism
\[
 \iota:B\longrightarrow\operatorname{Pic}^{\tau}_{X/R}
                                                        \tag{D4a.2}
\]
whose special fibre identifies $B_k$ isomorphically with
$\operatorname{Pic}^{\rm var}_{X_0/k}$ and whose special-fibre
map is its natural inclusion. In particular no abelian closed
subscheme of the relative Picard scheme can have the specified
reduced identity component on the closed fibre. The assertion
rules out existence, not merely a preferred construction.

Here is the proof. Put $Y_0=F_0/\langle P\rangle$ and write
$\phi:F_0\to Y_0$ for the étale degree-two isogeny. The product
Igusa calculation gives
\[
 \operatorname{Pic}^{\tau}_{X_0/k}
     \simeq Y_0^\vee\times E_0^\vee[2],
 \qquad \operatorname{Pic}^{\rm var}_{X_0/k}=Y_0^\vee.
\]
Under pullback to $A_0$, the latter inclusion becomes
\[
 f_0=(0,\phi^\vee):Y_0^\vee
                         \longrightarrow E_0^\vee\times F_0^\vee.
                                                        \tag{D4a.3}
\]
This also follows directly by pulling back a degree-zero line
bundle from the quotient map $X_0\to Y_0$.
If (D4a.2) existed, compose it with $\pi^*$ to the relative
Picard scheme of $A$. This composite lands in its open neutral
component $A^\vee$: it does so on the closed fibre, and an
Artinian nilpotent thickening has the same underlying topological
space, so the inverse image of the complementary open is empty.
We therefore obtain a homomorphism of ordinary abelian $R$-schemes
\[
 f:B\longrightarrow A^\vee
\]
lifting $f_0$. Both source and target are abelian schemes, so its
dual $f^\vee:A\to B^\vee$ is a homomorphism as well; no image,
kernel-flatness, or splitting assertion is needed.

Choose the ordinary étale Tate bases for $A_0$ and $A_0^\vee$
as in D4o, with indices $E,F$. Since $\ker\phi^\vee=\mu_2$,
its map on étale $2$-adic Tate modules is an isomorphism. Choose
the generator $x\in T_2Y_0^\vee(k)$ so that
\[
 T_2(f_0)(x)=e_F^\vee.
\]
The dual of (D4a.3) is $(0,\phi):E_0\times F_0\to Y_0$,
so
\[
 T_2(f_0^\vee)(e_E)=0.
\]
Katz's necessary and sufficient morphism-lifting criterion
\cite[Theorem 2.1(4), p.~149]{DGJ_Katz1981}, applied to
$f:B\to A^\vee$, would imply
\[
 q_B\bigl(x,T_2(f_0^\vee)e_E\bigr)
    =q_{A^\vee}\bigl(T_2(f_0)x,e_E\bigr).
\]
The left side equals one. By the transposition formula for dual
abelian schemes, Katz's Theorem~2.1(3), pp.~148--149, the right
side equals
\[
 q_{A^\vee}(e_F^\vee,e_E)=q_A(e_E,e_F^\vee)
      =q_{12}=1+\epsilon\ne1.
\]
This contradiction proves the assertion. It uses the exact
Serre--Tate lifting criterion on the special-fibre Tate modules,
not a claim that an arbitrary point lifts through a torsion
isogeny. It remains valid for every possible elliptic deformation
$B$ of $Y_0^\vee$: its first pairing entry is evaluated at zero
in the second slot, independently of its own parameter.

The leaf 57:311 asks ``Existence de'' $\operatorname{Pic}^{00}$
(the superscript is transcribed uncertainly as $00$ or $\infty$),
``dans le cas Artinien'' with ``local'' uncertain. Under the
compatible-abelian-subscheme interpretation specified above,
the answer is negative already on an Artinian base of length two. The full torsion Picard scheme
of this same example is flat. The two properties must not be
identified. The calculation of the matrix and the nonzero
obstruction in $R$ is certified in
\texttt{check\_artin\_picvar.py}; existence and the contradiction
are the geometric proof above.

The targeted independent audits by root and bce checked the étale Tate
isomorphism for $\phi^\vee$, the zero dual entry, the exact
transposition convention, and the existence of the quotient from
D4o. No gap was found. This explicit Serre--Tate proof is consistent
with Grothendieck's published first-order Igusa remark
\cite[Remark 3.4, p.~238]{DGJ_FGAVI1962}, not a claim of priority.
